\chapter{Cykl roczny}
�wi�ta cyklu rocznego cechuj� si� powtarzalno�ci� na przestrzeni lat - jest to tzw. �Ko�o obrz�dowe�. Poni�ej prezentujemy sugerowane, orientacyjne daty ich obchod�w. Warto mie� na uwadze, �e nasi przodkowie obchodzili �wi�ta cz�sto �na wyczucie� tak, aby odpowiada�y one cyklom przyrody. W efekcie daty te mog� si� na przestrzeni terytorium RP r�ni� nawet o 3 tygodnie. Warto te� organizowa� �wi�ta w soboty i niedziele, aby korzysta� z dni ustawowo wolnych od pracy.
\begin{itemize}
    \item \textbf{Zapusty} -- 6.01--15.02
    \item \textbf{Gromnica} -- 02.02--15.02
    \item \textbf{Welesowe} -- 03.02--11.02
    \item \textbf{Jare Gody} -- 21.03
    \item \textbf{Matki Ziemi} -- 22.04
    \item \textbf{Dziady Wiosenne}\footnote{Zobacz opis �wi�ta.} -- 02.05
    \item \textbf{Rusalny tydzie�} -- 3.05--10.05
    \item \textbf{Zielone �wi�tki / Maik / Stado} -- 11.05--08.06
    \item \textbf{Noc Kupa�y} -- 20--22.06
    \item \textbf{�wi�to Peruna} -- 15.07
    \item \textbf{Trzy spasy letnie i Leszy} -- 14--29.08
    \item \textbf{�wi�to Mokoszy} -- 15.08
    \item \textbf{�wi�to Plon�w} -- 22.09
    \item \textbf{Dziady Jesienne}\footnote{Zobacz opis �wi�ta.} -- 01.11
    \item \textbf{Spas Zimowy} -- 05.12
    \item \textbf{Szczodre Gody} -- 20.12--06.01
\end{itemize}